% macros.tex — notations du mémoire (ML distribué, Byzantin, code)
% Convention : commandes en français/anglais explicites, rien de spéculatif.

% --- Texte ---
\newcommand{\textcode}[1]{\texttt{#1}}

% --- Éditions colorées (relecture) : génère \arthuredit{}, \arthurreplace{}{},
% \arthurcomment{}, \arthurdelete{} ; option [suppress] pour tout masquer ---
\addauthor{arthur}{red}

% --- Maths de base ---
\newcommand{\vect}[1]{\boldsymbol{#1}} % vecteur / gradient
\newcommand{\norm}[1]{\lVert #1 \rVert} % norme euclidienne
\newcommand{\sumn}{\sum_{i=1}^{n}} % somme sur n workers

% --- Environnements mathématiques ---
% Requiert amsmath, amssymb et amsthm chargés avant cet input.
% Compteur partagé, numérotation par section : valable en article comme en report.
% Preuve et fin de preuve : environnement proof fourni par amsthm.
\theoremstyle{plain}
\newtheorem{theorem}{Théorème}[section]
\newtheorem{lemma}[theorem]{Lemme}
\newtheorem{proposition}[theorem]{Proposition}
\newtheorem{corollary}[theorem]{Corollaire}
\newtheorem{property}[theorem]{Propriété}
\theoremstyle{definition}
\newtheorem{definition}[theorem]{Définition}
\newtheorem{assumption}[theorem]{Hypothèse}
\theoremstyle{remark}
\newtheorem{remark}[theorem]{Remarque}
\newtheorem{example}[theorem]{Exemple}

% --- Opérateurs mathématiques ---
\DeclareMathOperator*{\argmin}{arg\,min}
\DeclareMathOperator*{\argmax}{arg\,max}

% --- Apprentissage distribué ---
\newcommand{\grad}{g} % gradient stochastique
\newcommand{\loss}{\mathcal{L}} % fonction objectif
\newcommand{\dataset}{\mathcal{D}} % jeu de données
\newcommand{\model}{w} % paramètres du modèle

% --- Byzantin ---
\newcommand{\nworkers}{n} % nombre total de workers
\newcommand{\nbyz}{f} % nombre de malicieux
\newcommand{\nselected}{m} % nombre de gradients sélectionnés par MultiKrum
\newcommand{\honest}{\mathcal{H}} % ensemble des honnêtes
\newcommand{\byzset}{\mathcal{B}} % ensemble des malhonnêtes

% --- Nom de la librairie (robuste aux signets PDF) ---
% \krum (minuscules, monospace) désigne la règle d'agrégation, comme les
% autres macros d'agrégateurs ci-dessous ; ne pas confondre.
\DeclareRobustCommand{\krumlib}{\texorpdfstring{\textsc{Krum}}{Krum}}

% --- Agrégateurs (noms d'algorithmes) ---
\newcommand{\krum}{\textcode{Krum}}
\newcommand{\mkrum}{\textcode{MultiKrum}}
\newcommand{\monna}{\textcode{MoNNA}}
\newcommand{\bulyan}{\textcode{Bulyan}}
\newcommand{\tmean}{\textcode{TrimmedMean}}
\newcommand{\average}{\textcode{Average}}
\newcommand{\median}{\textcode{Median}}
\newcommand{\geomed}{\textcode{GeoMed}}
\newcommand{\medoid}{\textcode{Medoid}}
\newcommand{\aksel}{\textcode{Aksel}}
\newcommand{\brute}{\textcode{Brute}}
\newcommand{\nna}{\textcode{NNA}}
\newcommand{\rfa}{\textcode{RFA}}
% Travaux futurs (ch. 7) : règle 1-bit, visée mais pas encore implémentée.
\newcommand{\signsgd}{\textcode{signSGD}}

% --- Attaques (noms courts, sans le suffixe Attack des classes) ---
\newcommand{\signflip}{\textcode{SignFlip}}
\newcommand{\alie}{\textcode{ALIE}}
\newcommand{\gaussian}{\textcode{Gaussian}}
\newcommand{\fullgradientnegation}{\textcode{FullGradientNegation}}
\newcommand{\smallperturbation}{\textcode{SmallPerturbation}}

% --- Code inline ---
\newcommand{\code}[1]{\texttt{#1}}
% Coupure de ligne autorisée mais invisible, pour les longs identifiants
% CamelCase dans les colonnes étroites : \code{Full\coupure{}Gradient}.
\newcommand{\coupure}{\allowbreak{}}